\documentclass[11pt]{article}
\usepackage[a4paper,margin=26mm]{geometry}
\usepackage[T1]{fontenc}
\usepackage{lmodern,microtype,amsmath,amssymb,amsthm,mathtools,booktabs}
\usepackage{xcolor}
\usepackage[colorlinks=true,linkcolor=blue!45!black,citecolor=blue!45!black,urlcolor=blue!45!black]{hyperref}
\usepackage{fancyhdr}
\pagestyle{fancy}\fancyhf{}\fancyhead[L]{\small Sharp graph and entire-function inequalities}
\fancyhead[R]{\small Research note}\fancyfoot[C]{\thepage}
\setlength{\headheight}{14pt}
\newtheorem{theorem}{Theorem}[section]
\newtheorem{lemma}[theorem]{Lemma}
\newtheorem{corollary}[theorem]{Corollary}
\newtheorem{proposition}[theorem]{Proposition}
\theoremstyle{remark}\newtheorem{remark}[theorem]{Remark}
\newcommand{\LP}{\mathcal{LP}}
\newcommand{\R}{\mathbb R}
\newcommand{\C}{\mathbb C}
\newcommand{\W}{\mathcal W}
\DeclareMathOperator{\mesh}{mesh}
\title{\textbf{Sharp Line-Graph Growth and a Gaussian\newline Discrete Laguerre Inequality}}
\author{Research note prepared with ChatGPT for Keio}
\date{15 November 2026}
\begin{document}
\maketitle
\begin{abstract}
Two self-contained global inequalities are proved. For a finite connected
simple graph $G$ of order $n$ and minimum degree $\delta\ge2$, the modified
Randi\'c index of its line graph is at least $\delta n/4$. Consequently
$R(L(G))\ge(\delta/2)R(G)$, with equality exactly for regular graphs.
The factor is optimal and implies the statement of Zhang--Wu Conjecture~5
(2022). An exact nonnegative defect formula and a leaf-support extension
are also obtained. For a nonzero Laguerre--P\'olya entire function $\varphi$
with zero spacing at least $h$, and Gaussian factor $e^{-az^2}$, we prove
$D_h(\varphi;x)\ge4(1-e^{-2ah^2})\varphi(x)^2$ for every real $x$.
Both the Gaussian coefficient and the spacing threshold are sharp.
This includes the statement of Chasse--Csordas Conjecture~3.5 (2012).
The proofs do not use unverified corpus theorems or computational evidence
as assumptions. Historical priority is not asserted; external mathematical
review and a comprehensive literature check remain appropriate.
\end{abstract}

\section{The two problems and the scope of the claims}

Zhang and Wu \cite[Conjecture~5]{ZW} ask whether the Randi\'c index cannot
decrease under the line-graph operation when the minimum degree is at
least two, with equality only for cycles. Liang and Wu
\cite[Theorem~1.4]{LW} prove this under the additional restriction
$\Delta\le\delta^2-2\delta+2$. Theorem~\ref{thm:graph} below removes that
restriction and gives the optimal multiplier $\delta/2$.

Chasse and Csordas \cite[Conjecture~3.5]{CC} ask for a discrete Laguerre
inequality for the full Laguerre--P\'olya class under a unit zero-spacing
condition. Their Theorem~3.4 handles order below two and minimal type at
order two. Theorem~\ref{thm:LP} below includes a nonzero Gaussian factor
and gives a sharp positive correction. This is a different statement
from their finite-difference level-set converse, Conjecture~2.22.

These are independently checkable derivations of the stated mathematical
claims. A targeted search of primary sources did not locate a later
general settlement of either target, but that is not a proof of novelty.
Both proofs are self-contained: the graph argument uses an exact allocation
over incident edge pairs, and the entire-function argument combines zero
spacing with logarithmic concavity.

\section{An optimal line-graph multiplier}

All graphs in this note are finite and simple. For a graph $H$, write
\[
 R(H)=\sum_{uv\in E(H)}\frac1{\sqrt{d_H(u)d_H(v)}},\qquad
 M(H)=\sum_{uv\in E(H)}\frac1{\max\{d_H(u),d_H(v)\}}.
\]
The latter is often called the modified Randi\'c index. Write
$\mathsf H(H)=\sum_{uv\in E(H)}2/(d_H(u)+d_H(v))$ for the harmonic index.
Termwise,
\begin{equation}\label{eq:index-chain}
 M(H)\le\mathsf H(H)\le R(H).
\end{equation}
The vertices of $L(G)$ are the edges of $G$; distinct edges of $G$ are
adjacent in $L(G)$ precisely when they share an endpoint.

\begin{theorem}[Sharp minimum-degree multiplier]\label{thm:graph}
Let $G$ be connected, of order $n$, with minimum degree $\delta\ge2$.
Then
\begin{equation}\label{eq:main-graph}
 M(L(G))\ge\frac{\delta n}{4},\qquad
 R(L(G))\ge\frac{\delta}{2}R(G).
\end{equation}
Equality in the second inequality holds if and only if $G$ is regular.
The coefficient $\delta/2$ is best possible for every integer $\delta\ge2$.
In particular, $R(L(G))\ge R(G)$, with equality if and only if $G$ is a cycle.
\end{theorem}

\begin{proof}
Write $d_v=d_G(v)$ and $m=|E(G)|$. The line-graph degree of an edge $uv$ is
\[
 D_{uv}=d_u+d_v-2.
\]
For a fixed vertex $v$, put $x_u=D_{uv}^{-1}$ for $u\in N(v)$, and
$c_v=(d_v+\delta-2)^{-1}$. Then $0<x_u\le c_v$.
Each edge of the line graph is owned by a unique unordered pair of
distinct edges meeting at a vertex $v$. Consequently
\begin{align*}
 M(L(G))
 &=\sum_v\sum_{\{u,w\}\subset N(v)}\min(x_u,x_w)\\
 &\ge\sum_v\left[(d_v-1)\sum_{u\in N(v)}\frac1{D_{uv}}
             -\binom{d_v}{2}c_v\right].
\end{align*}
Here $\min(x,y)=x+y-\max(x,y)\ge x+y-c_v$ was used.
The two endpoint contributions of an original edge $uv$ in the first
sum are exactly
\[
 \frac{d_u-1}{D_{uv}}+\frac{d_v-1}{D_{uv}}=1.
\]
Thus that sum is $m$. Using $\sum_vd_v=2m$ gives
\begin{align}
 M(L(G))
 &\ge m-\frac12\sum_v\frac{d_v(d_v-1)}{d_v+\delta-2}\notag\\
 &=\frac{\delta-1}{2}\sum_v\frac{d_v}{d_v+\delta-2}
 \ge\frac{\delta n}{4}.\label{eq:graph-middle}
\end{align}
For the last step, $d_v\ge\delta$ and
\[
 \frac{d_v}{d_v+\delta-2}\ge\frac{\delta}{2\delta-2};
\]
this is an equality for all $d_v$ when $\delta=2$ and follows by
cross-multiplication when $\delta>2$.

For every graph without isolated vertices,
\begin{equation}\label{eq:Rupper}
 R(G)\le\frac12\sum_{uv\in E(G)}\left(\frac1{d_u}+\frac1{d_v}\right)
 =\frac n2.
\end{equation}
Combining \eqref{eq:index-chain}, \eqref{eq:graph-middle}, and
\eqref{eq:Rupper} proves the second inequality.

If equality holds in that inequality, equality must hold in
\eqref{eq:Rupper}. Hence $d_u=d_v$ on every edge, and connectedness makes
$G$ regular. Conversely, for a $d$-regular graph, $R(G)=n/2$, while
$L(G)$ is $(2d-2)$-regular on $nd/2$ vertices, so $R(L(G))=nd/4$.
This also proves optimality. Equality in $R(L(G))\ge R(G)$ additionally
forces $d=2$, which characterizes connected cycles.
\end{proof}

\subsection{An exact certificate of the gap}

Let $\W$ be the set of triples $(v,\{u,w\})$ with $u,w$ distinct neighbors
of $v$, and put $t_{u,w}=\min(d_u,d_w)$. The preceding proof gives the exact
identity
\begin{align}
 M(L(G))-\frac{\delta n}{4}
 ={}&\frac{\delta-2}{4}\sum_v\frac{d_v-\delta}{d_v+\delta-2}\notag\\
 &+\sum_{(v,\{u,w\})\in\W}
 \frac{t_{u,w}-\delta}
 {(d_v+\delta-2)(d_v+t_{u,w}-2)}.\label{eq:exact-delta}
\end{align}
Every summand is nonnegative. Also,
\[
 \frac n2-R(G)
 =\frac12\sum_{uv\in E(G)}
       \left(\frac1{\sqrt{d_u}}-\frac1{\sqrt{d_v}}\right)^2.
\]
Therefore $R(L(G))-(\delta/2)R(G)$ is exactly the sum of the
nonnegative right side of \eqref{eq:exact-delta}, the termwise gap
$R(L(G))-M(L(G))$, and
\[
 \frac\delta4\sum_{uv\in E(G)}
       \left(\frac1{\sqrt{d_u}}-\frac1{\sqrt{d_v}}\right)^2.
\]
This explicitly records the sources of strictness.

\begin{corollary}[Iterated line graphs]
Under the hypotheses of Theorem~\ref{thm:graph}, for every integer $k\ge1$,
\[
 R(L^k(G))\ge R(G)\prod_{j=0}^{k-1}
 \left(1+(\delta-2)2^{j-1}\right).
\]
Equality holds for every $\delta$-regular graph. In particular, the Randi\'c
indices of all successive line graphs form a nondecreasing sequence.
\end{corollary}
\begin{proof}
The degree formula gives $\delta(L(H))\ge2\delta(H)-2$.
Inductively, $\delta(L^j(G))\ge2+(\delta-2)2^j$. Apply
Theorem~\ref{thm:graph} successively. For a regular initial graph all
intermediate graphs are regular, and every bound is an equality.
\end{proof}

\section{A leaf-support bound without a minimum-degree hypothesis}

\begin{theorem}\label{thm:support}
Let $G$ be connected of order $n\ge3$, and let $s$ be the number of
vertices adjacent to at least one leaf. Then
\begin{equation}\label{eq:support}
 M(L(G))\ge\frac{n-s}{2}\ge\frac n4.
\end{equation}
Equality in the last bound $M(L(G))\ge n/4$ holds exactly for $G=P_4$.
In particular, for all such $G$,
\[
 \mathsf H(L(G))>\frac n4,\qquad R(L(G))>\frac n4.
\]
\end{theorem}

\begin{proof}
All $D_{uv}>0$ because a connected graph of order at least three has no
component $K_2$. The same pair identity, now adding and subtracting
$1/d_v$, gives
\begin{equation}\label{eq:all-graph-identity}
 M(L(G))=\frac n2+
 \sum_{(v,\{u,w\})\in\W}
 \left(\frac1{d_v}-\frac1{d_v+\min(d_u,d_w)-2}\right).
\end{equation}
This holds also at degree-one vertices, which have no pairs.
Let $\ell_v$ be the number of leaf neighbors of $v$. At a vertex of degree
$d_v\ge2$, precisely $\ell_v(2d_v-\ell_v-1)/2$ pairs contain a leaf.
Each contributes $-1/[d_v(d_v-1)]$. The other pairs contribute nonnegatively.
The total negative contribution at $v$ is therefore at most
\[
 P_v=\frac{\ell_v(2d_v-\ell_v-1)}{2d_v(d_v-1)}\le\frac12,
\]
where the last inequality follows from
\[
 d_v(d_v-1)-\ell_v(2d_v-\ell_v-1)
 =(d_v-\ell_v)(d_v-\ell_v-1)\ge0.
\]
Only the $s$ supports contribute such a penalty. This proves the first
bound. Each support is a nonleaf and has a distinct leaf assigned to it,
so $2s\le n$.

If $M(L(G))=n/4$, equality forces $s=n/2$. There are then exactly $s$
leaves and $s$ nonleaves; every nonleaf has exactly one leaf neighbor.
Now $P_v=1/d_v$ on each support, so equality forces every nonleaf to
have degree two. The graph induced by the nonleaves is connected and
1-regular, hence is $K_2$. Thus $G=P_4$. Conversely,
$M(L(P_4))=M(P_3)=1$. Finally,
$\mathsf H(P_3)=4/3>1$ and $R(P_3)=\sqrt2>1$, proving the strict
consequences even in that exceptional case.
\end{proof}

The tree-only harmonic consequence had already been studied by 2023,
as recorded in \cite{ZZ}; it is not presented as a new historical claim.
The proof here gives the support-sensitive bound for every connected
simple graph, including graphs with cycles.

\section{A sharp Gaussian discrete Laguerre inequality}

A nonzero real entire function belongs to the Laguerre--P\'olya class if
it has a representation
\begin{equation}\label{eq:LP-product}
 \varphi(z)=Cz^k e^{-az^2+bz}
 \prod_j\left(1-\frac z{r_j}\right)e^{z/r_j},
\end{equation}
where $C\ne0$, $b,r_j\in\R$, $a\ge0$, $r_j\ne0$, and
$\sum_jr_j^{-2}<\infty$. The product may be finite or empty.
Zeros are counted \emph{with multiplicity}. We write $\mesh(\varphi)\ge h$
when distinct entries in that multiset are separated by at least $h$;
the condition is vacuous for at most one zero. In particular, $h>0$
excludes all multiple zeros.

Set
\begin{align}
 D_h(\varphi;x)={}&[\varphi(x+h)-\varphi(x-h)]^2\notag\\
 &-4\varphi(x)[\varphi(x+h)-2\varphi(x)+\varphi(x-h)].\label{eq:Dh}
\end{align}

\begin{theorem}[Sharp Gaussian correction]\label{thm:LP}
Suppose $\varphi\in\LP$ is nonzero and $\mesh(\varphi)\ge h>0$.
If $a$ is its Gaussian coefficient in \eqref{eq:LP-product}, then for
every real $x$,
\begin{equation}\label{eq:LP-main}
 D_h(\varphi;x)\ge4(1-e^{-2ah^2})\varphi(x)^2\ge0.
\end{equation}
For each $a\ge0$ and $h>0$, the first coefficient is optimal over this class.
\end{theorem}

\begin{proof}
If $\varphi(x)=0$, \eqref{eq:LP-main} is the nonnegativity of a square.
Otherwise define
\[
 A=\frac{\varphi(x+h)}{\varphi(x)},\qquad
 B=\frac{\varphi(x-h)}{\varphi(x)},\qquad
 Q=\frac{D_h(\varphi;x)}{\varphi(x)^2}.
\]
There are two useful identical expressions:
\begin{equation}\label{eq:square-split}
 Q=(A-B)^2-4(A+B-2)=(A+B-2)^2+4(1-AB).
\end{equation}
If $AB\le0$, the second expression is at least $4$.
If $A,B<0$, the first expression is greater than $8$.
Both bounds exceed or equal $4(1-e^{-2ah^2})$, which is less than $4$.

It remains to treat $A,B>0$. The three values
$\varphi(x-h),\varphi(x),\varphi(x+h)$ have the same sign.
Each open interval $(x-h,x)$ and $(x,x+h)$ contains at most one zero
by the spacing assumption. Every zero is simple, so one zero would
reverse the endpoint signs. Neither interval can therefore contain a
zero. Their endpoints are also nonzeros.

On the zero-free interval $[x-h,x+h]$, let $g(t)=\log|\varphi(t)|$.
Logarithmic differentiation of \eqref{eq:LP-product} gives
\begin{equation}\label{eq:log-curvature}
 g''(t)=-2a-\sum_{r:\,\varphi(r)=0}\frac1{(t-r)^2}\le-2a.
\end{equation}
The zero at the origin, if present, is included in the sum. The series
converges locally uniformly off the zeros because $\sum_jr_j^{-2}<\infty$.
Thus $g(t)+at^2$ is concave on this interval, and its midpoint inequality is
\[
 g(x+h)+g(x-h)-2g(x)\le-2ah^2.
\]
Exponentiating yields $AB\le e^{-2ah^2}$. Substituting into
\eqref{eq:square-split} proves \eqref{eq:LP-main} in the remaining case.

For sharpness, take the zero-free function $\varphi(x)=e^{-ax^2+bx}$.
A direct substitution gives
\begin{align}
 \frac{D_h(\varphi;x)}{\varphi(x)^2}
 ={}&\left[2e^{-ah^2}\cosh\bigl(h(b-2ax)\bigr)-2\right]^2\notag\\
 &+4(1-e^{-2ah^2}).\label{eq:Gaussian-sharp}
\end{align}
For $a>0$, a real $x$ can be chosen so that
$\cosh(h(b-2ax))=e^{ah^2}$. For $a=0$, a nonzero constant function
attains equality. The coefficient cannot be increased.
\end{proof}

\begin{corollary}[The stated Chasse--Csordas conjecture]
If $\varphi\in\LP$ and its zeros, counted with multiplicity, have spacing
at least one, then $D_1(\varphi;x)\ge0$ for every real $x$.
\end{corollary}

If $\varphi$ has at least one real zero, \eqref{eq:log-curvature} is
strictly less than $-2a$ on every zero-free interval. The strengthened
inequality is consequently strict at every $x$ where $\varphi(x)\ne0$.
At a zero $x$, equality is equivalent to
$\varphi(x+h)=\varphi(x-h)$; the spacing and simplicity imply this occurs
exactly when both neighboring values are zero.

\section{The zero-spacing threshold is also sharp}

\begin{proposition}\label{prop:mesh-sharp}
For every $h>0$ and every $0<d<h$, there is a function
$\varphi\in\LP$ with exactly two simple zeros separated by $d$ for which
$D_h(\varphi;0)<0$.
\end{proposition}

\begin{proof}
Let
\[
 r_1=\frac{h-d}{2},\qquad r_2=\frac{h+d}{2},\qquad
 P(x)=(x-r_1)(x-r_2)=x^2-hx+q,
\]
where $q=(h^2-d^2)/4>0$. Put
\[
 K=1+\frac{2h^2}{q}>1,\qquad
 b=\frac{\log K}{2h},\qquad \varphi(x)=e^{bx}P(x).
\]
This is a Laguerre--P\'olya function with mesh exactly $d$.
Since $P(0)=P(h)=q$ and $P(-h)=qK$,
\[
 \frac{\varphi(h)}{\varphi(0)}
 =\frac{\varphi(-h)}{\varphi(0)}=\sqrt K.
\]
The first expression in \eqref{eq:square-split} therefore gives
\[
 D_h(\varphi;0)=8q^2(1-\sqrt K)<0.
\]
Thus no smaller uniform spacing threshold can replace $h$.
\end{proof}

The multiplicity convention is essential as well. If one measured only
the distance between distinct zero locations, the polynomial
$\varphi(x)=(x^2-1/4)^2$ would have apparent spacing one, but
$D_1(\varphi;0)=-1/4$. Its actual mesh, counting multiplicity, is zero.

\section{Verification and mathematical status}

The proofs above are symbolic and apply to all objects satisfying their
hypotheses. Exact rational graph enumeration was used only as a check of
the identities and equality cases. All connected labelled simple graphs
of orders three through six were examined:
\begin{center}
\begin{tabular}{@{}rrr@{}}\toprule
Order & Connected graphs & Minimum degree at least two\\\midrule
3 & 4 & 1\\
4 & 38 & 10\\
5 & 728 & 253\\
6 & 26\,704 & 12\,058\\\midrule
Total & 27\,474 & 12\,322\\\bottomrule
\end{tabular}
\end{center}
The exact leaf decomposition, support bound, minimum-degree bound and
their rational defect identities passed these checks. The only labelled
graphs attaining $M(L(G))=n/4$ were the twelve labelings of $P_4$.

For the entire-function result, the proof explicitly covers zeros at
the center or endpoints, both same-sign and opposite-sign ratios, finite
and infinite canonical products, and the Gaussian factor. The coefficient
and mesh sharpness claims have explicit exact witnesses, rather than
numerical minimization.

Theorems~\ref{thm:graph} and \ref{thm:LP} establish the respective statements
as formulated in the cited original papers. They are research proofs
prepared in this session, not externally refereed publications. Their
priority is separate from their mathematical content.

\begin{thebibliography}{9}
\bibitem{ZW}
J. Zhang and B. Wu,
\emph{Randi\'c Index of a Line Graph},
Axioms \textbf{11}(5) (2022), 210.
\href{https://doi.org/10.3390/axioms11050210}{doi:10.3390/axioms11050210}.

\bibitem{LW}
Y. Liang and B. Wu,
\emph{General Randi\'c indices of a graph and its line graph},
Open Mathematics \textbf{21}(1) (2023), 20220611.
\href{https://doi.org/10.1515/math-2022-0611}{doi:10.1515/math-2022-0611}.

\bibitem{ZZ}
Z. Zou and L. Zhong,
\emph{A lower bound on the modified Randi\'c index of line graphs},
Discrete Applied Mathematics \textbf{345} (2024), 136--140.
\href{https://doi.org/10.1016/j.dam.2023.11.050}{doi:10.1016/j.dam.2023.11.050}.

\bibitem{CC}
M. Chasse and G. Csordas,
\emph{Discrete analogues of the Laguerre inequalities and a conjecture of
I. Krasikov},
Annales des sciences math\'ematiques du Qu\'ebec \textbf{36}(1) (2012), 77--91.
\href{https://www.labmath.uqam.ca/~annales/volumes/36-1/PDF/077-091.pdf}{Original article}.
\end{thebibliography}
\end{document}
